\input{preamble.tex}
\begin{document}

\begin{titlepage}
\vspace*{18mm}
{\sffamily\small ПРОЕКТНЫЙ ДОКУМЕНТ · ЧЕРНОВИК 0.2 · 10 НОЯБРЯ 2026}\par
\vspace{8mm}
{\sffamily\bfseries\fontsize{26}{30}\selectfont Тенгварская пишущая\\[2pt] машинка ТПМ-2}\par
\vspace{4mm}
{\large Механическое исполнение на челноке Хаммонда:\\ подвижный печатающий узел и память на трёх секторах}\par
\vspace{14mm}
\begin{center}
\setlength\tgem{13.5mm}
\begin{tikzpicture}[x=\tgu,y=\tgu]
\input{gen/fig_trace_q.tex}
\end{tikzpicture}
\end{center}
\vspace{6mm}
\noindent\begin{minipage}{0.86\textwidth}\small
\textbf{О чём документ.} Редакция 0.2 закрепляет конструкцию, выбранную по двум критериям — простота в эксплуатации и точность механизма. ТПМ-2 — ручная машинка без мотора. Все знаки отлиты на одном сменном челноке; бумагу на челнок бьёт молоток с постоянной силой пружины. Продвигающие литеры печатаются от левого края клетки, после удара каретка проходит ширину литеры. Для мёртвого знака весь печатающий узел (челнок с молотком на общей раме) отходит назад на величину, которую запомнила предыдущая буква, — на $b$ для надстрочного знака, на $c$ для подстрочного, — и после удара возвращается на нулевой упор. Оттиск совпадает с оттиском ТПМ-1 знак в знак (проверено моделью). Рычажное исполнение с салазками (редакция 0.1) и другие варианты вынесены в приложение.
\end{minipage}
\vfill
{\small User · продолжение проектов ТПМ-1 и ТПМ-2 (ред. 0.1).}
\end{titlepage}

\tableofcontents
\newpage

%=====================================================================
\section{Выбор конструкции}

\subsection{Критерии}

\begin{description}[style=nextline,leftmargin=0pt]
  \item[Простота в эксплуатации.] Наборщику не нужно знать устройство машинки: удобная раскладка, одно правило порядка ударов, исправление одним нажатием, ровный оттиск при любой силе нажатия, смена шрифта или режима без инструмента.
  \item[Точность механизма.] Положение каждого знака задаётся как можно меньшим числом звеньев, каждое звено доводится жёстким упором или фиксатором до целой единицы, и регулировка, от которой зависит посадка всех знаков, одна.
\end{description}

\subsection{Решения}

\begin{center}\small
\begin{tabularx}{\linewidth}{@{}>{\raggedright}p{33mm}>{\raggedright}p{34mm}X@{}}
\toprule
Вопрос & Решение & Почему \\
\midrule
печатающий механизм & челнок Хаммонда & один литероноситель и одна точка удара: выравнивание задано конструкцией; молоток бьёт с постоянной силой; челнок сменный \\
как попасть на ось & отвод всего печатающего узла & одно подвижное звено с нулевым упором вместо двух салазок в сегменте \\
нижняя ось & сохранить & без неё подстрочный знак под большинством основных тенгв уходит на 0,4~мм, под hyarmen — на 1,2~мм; стоит только второго сектора памяти \\
«назад на литеру» & включить & исправление одним нажатием вместо $w$ ударов забоя (до 16); третий сектор памяти \\
фиксаторы & на всех секторах и на упоре спуска & прощают неточность зазоров меньше половины ступени \\
контроль хода и очерёдность & сохранить & исключают недожатую клавишу и перекрытие нажатий ценой скорости \\
ширины & без сведения к классам & меньше уровней подачи — проще упор, но промежутки поплывут почти на единицу \\
раскладка & как у ТПМ-1 & челнок снимает ограничение «клавиша либо мёртвая, либо продвигающая»: Shift + гласная снова даёт долгую гласную на ára \\
\bottomrule
\end{tabularx}
\end{center}

\paragraph{Что не изменилось.} Литеры, их размеры и оси измерены на шрифте Alcarin Tengwar и пересчитаны в единицы $u = \mathrm{em}/16$, как в ТПМ-1. Кегль: em~=~1/4~дюйма, единица $u$~=~1/64~дюйма~$\approx$~0,4~мм. Порядок ударов: сначала тенгва, потом её знаки, в любом режиме.

%=====================================================================
\section{Предшественники}

\begin{description}[style=nextline,leftmargin=0pt]
  \item[Хаммонд (1880-е).] Машинка с литероносителем вместо рычагов. Литеры отлиты одним цельным челноком, который поворачивается на ободе кольцевой наковальни; нужная литера встаёт перед молотком, и тот под действием пружины бьёт бумагу о литеру. Задача такой схемы — идеальное выравнивание и равномерный оттиск. Сила молотка постоянна, поэтому оттиск ровный, как бы сильно или слабо ни нажимали клавиши. Челнок сменный, отсюда сотни шрифтов и языков. Наковальня и молоток сидят на U-образной раме. Недостаток — заправка нового листа требует лишних усилий.
  \item[Varityper.] Наследник Хаммонда с пропорциональной подачей: колесо спуска проходит две, три или четыре единицы в зависимости от нажатой клавиши, и спуск связан с молотком так, чтобы каретка не трогалась раньше удара.
  \item[Комптометр (1887).] Механизм работает от пальца: клавиша сама задаёт точку останова сектора и сама его поворачивает. Ход клавиши длинный; в поздних моделях машина блокировалась при недожатой клавише.
\end{description}

%=====================================================================
\section{Принцип}

\subsection{Геометрия}

Пусть точка удара (ось молотка при нулевом положении узла) находится против точки $P$ на бумаге.

\begin{enumerate}
  \item Продвигающая литера отлита на челноке так, что левый край её клетки приходится на ось молотка. Удар кладёт её клетку от $P$; ось техты оказывается в $P + a$.
  \item При отпускании клавиши каретка проходит $w$; против оси молотка теперь $P + w$.
  \item Для надстрочного знака узел отходит назад на $b$: удар в $P + w - b = P + a$, на ось буквы.
  \item Для подстрочного знака узел отходит на $c = b - d$: удар в $P + w - c = P + a + d$, на нижнюю ось.
\end{enumerate}

\begin{figure}[h]
\centering
\setlength\tgem{48mm}
\begin{tikzpicture}[x=\tgu,y=\tgu]
\input{gen/fig_geometry.tex}
\end{tikzpicture}
\caption{tinco, акут и две точки снизу. Красным — ось молотка при ударе литеры (1) и после подачи (2). Удар надстрочного знака (3) — с отводом узла на $b$, подстрочного (4) — на $c$. Размеры в $u = \mathrm{em}/16$.}
\label{fig:geom}
\end{figure}

\subsection{Код литеры}

\begin{center}\small
\begin{tabular}{@{}clll@{}}
\toprule
& Что задаёт & Диапазон & Куда идёт \\
\midrule
$w$ & подачу каретки после удара & 4--16\,u (12 уровней) & переставной упор спуска; сектор W \\
$b$ & отвод узла для надстрочных знаков & 1--8\,u & сектор B \\
$c = b - d$ & отвод узла для подстрочных знаков & 1--9\,u & сектор C \\
\bottomrule
\end{tabular}
\end{center}

У всех литер $b < w$ и $c < w$: отвод никогда не уводит узел за левый край клетки. Мёртвые знаки несут только признак «над / под». Пробелы задают $w = 4$ или $6$ и обнуляют всю память.

\subsection{Проверка на модели}

Модель ТПМ-2 (\code{simulate\_mech} в \code{tools/gen\_shuttle.py}) прогоняется на обоих примерах (\S\ref{sec:examples}) вместе с моделью ТПМ-1; скрипт требует, чтобы каждая литера и каждый знак легли в одну и ту же точку бумаги, и останавливается при любом расхождении. Расхождений нет.

%=====================================================================
\section{Общая схема}

\begin{figure}[h]
\centering
\begin{tikzpicture}[x=1mm,y=1mm]
% paper on carriage
\fill[gray!8] (0,70) rectangle (160,76);
\draw[thinmech] (0,70) -- (160,70) (0,76) -- (160,76);
\node[lbl] at (22,73) {лист на каретке};
\draw[arr] (150,80) -- node[lbl,above]{подача $w$} (125,80);
% U-frame on slide
\draw[mech,fill=adv!60,rounded corners=1pt] (60,56) rectangle (100,62);
\node[lbl] at (80,59) {наковальня и челнок};
\draw[mech] (62,56) -- (62,46) -- (98,46) -- (98,56);
\draw[mech] (98,46) -- (110,46) -- (110,86) -- (86,86);
\draw[mech,fill=gray!20] (80,84) rectangle (86,88); \node[lbl,anchor=south] at (83,88.5) {молоток};
\draw[arr] (83,84) -- (83,77);
\node[lbl,anchor=west,align=left] at (111,66) {U-образная\\рама};
\fill[axisred] (83,70) circle (0.9);
% slide
\draw[mech,fill=gray!12] (40,40) rectangle (130,44);
\node[lbl,anchor=north] at (85,39.5) {ползун: ход назад до 9\,u (3,6 мм), параллельно листу};
\draw[thinmech] (80,44) -- (80,46);
% stops
\draw[mech,fill=ink!70] (132,38) rectangle (135,48); \node[lbl,anchor=west,align=left] at (136,43) {нулевой\\упор};
\draw[mech,fill=sledA!60] (34,46) rectangle (37,50); \draw[mech,fill=sledB!60] (30,38) rectangle (33,42);
\node[lbl,anchor=east,text=sledA] at (29,50) {упор $b$};
\node[lbl,anchor=east,text=sledB] at (29,37) {упор $c$};
\draw[{Latex[length=1.6mm]}-,thick,draw=accent] (44,51) -- ++(14,0) node[lbl,right,text=accent]{отвод};
% spring
\draw[decorate,decoration={coil,aspect=0.5,segment length=1.4mm,amplitude=1.2mm},thinmech] (130,42) -- (131.5,42);
% sectors, bails, escapement
\draw[mech,fill=white,rounded corners=1pt] (2,6) rectangle (52,22);
\node[lbl,align=center] at (27,14) {сектора памяти B, C, W\\с фиксаторами и защёлками};
\draw[sledA,thick] (20,22) -- (35.5,46); \draw[sledB,thick] (30,22) -- (31.5,38);
\draw[mech,fill=white,rounded corners=1pt] (58,6) rectangle (110,22);
\node[lbl,align=center] at (84,14) {кодовые пластины и мерные\\планки R, W, B, C, О, N};
\draw[mech,fill=white,rounded corners=1pt] (116,6) rectangle (158,22);
\node[lbl,align=center] at (137,14) {спуск: колесо, собачка,\\переставной упор $w$};
\draw[arr] (58,14) -- (52,14); \draw[arr] (110,14) -- (116,14);
\draw[thinmech] (150,22) -- (150,70);
\node[lbl,anchor=west,align=left] at (151,30) {к рейке\\каретки};
\draw[arr] (84,22) -- (84,40) node[pos=0.5,lbl,right]{О};
\end{tikzpicture}
\caption{ТПМ-2, вид сверху (условно). Челнок и молоток сидят на одной U-образной раме, рама — на ползуне. В покое ползун прижат пружиной к нулевому упору. Для мёртвого знака планка О отводит его назад до упора $b$ или $c$, который поставили сектора памяти; выбор упора делает признак N. Красная точка — ось молотка.}
\label{fig:top}
\end{figure}

\paragraph{Из чего состоит.} Остов машинки Хаммонда (челнок, наковальня, молоток, каретка с листом и пружинным двигателем, колесо спуска), к которому добавлены:
\begin{itemize}
  \item \textbf{ползун} под U-образной рамой с нулевым упором и пружиной возврата;
  \item \textbf{кодовые пластины} на всех 44 клавишах — по два яруса (регистра);
  \item \textbf{шесть мерных планок}: R (сброс памяти), W, B, C (замер) для продвигающих литер, О (отвод) и N (выбор упора) для мёртвых знаков;
  \item \textbf{переставной упор} спуска, который ставит планка W;
  \item \textbf{три сектора памяти} B, C, W с фиксаторами и защёлками;
  \item \textbf{блокировки} полного хода и очерёдности.
\end{itemize}

%=====================================================================
\section{Узлы}

\subsection{Печатающий узел}

\paragraph{Челнок.} Все 86 литер на одном цельном челноке: 4 ряда по 22 позиции (рис.~\ref{fig:shuttle}). Ряд выбирают две вещи: регистр (нижний / верхний) и половина клавиатуры (два задних ряда клавиш / два передних). Клавиши передних рядов несут на пластине выступ, который поднимает челнок на ряд; клавиша регистра — на два. Позицию в ряду задаёт поворот челнока, как у Хаммонда. Шаг позиции — 19\,u (самая широкая клетка 16\,u плюс 3\,u промежутка), так что развёртка ряда — около 166~мм.

\paragraph{Как литеры размещены на челноке.} Продвигающая литера отлита так, чтобы левый край её клетки при повороте челнока на свою позицию вставал против оси молотка. Мёртвый знак отлит так, чтобы против оси молотка вставал его якорь. Это единственное различие в отливке; на развёртке оно видно по положению красной линии.

\paragraph{Молоток.} Взводится ходом клавиши и спускается в конце хода; сила удара задана пружиной, а не пальцем. Это и даёт ровный оттиск тонких техт и широких двойных дуг.

\paragraph{Рама и ползун.} Челнок с наковальней и молоток сидят на одной U-образной раме, так что при отводе они смещаются вместе и ось удара переносится без перекоса. Рама стоит на ползуне с направляющими, параллельными листу. Ход ползуна — до 9\,u ($\approx$ 3,6~мм).

\begin{figure}[h]
\centering
\setlength\tgem{6.2mm}
\begin{tikzpicture}[x=\tgu,y=\tgu]
\input{gen/fig_shuttle.tex}
\end{tikzpicture}
\caption{Развёртка челнока: 4 ряда по 22 позиции. Серое поле — клетка продвигающей литеры, зелёные и коричневые позиции — надстрочные и подстрочные знаки. Красная линия — то, что встаёт против оси молотка: левый край клетки или якорь знака.}
\label{fig:shuttle}
\end{figure}

\subsection{Кодовая пластина и мерные планки}

Под клавишей висит кодовая пластина с выступами над мерными планками — длинными качающимися рамками поперёк клавиатуры. Код задан \emph{зазором}: при полном ходе $T$ выступ с зазором $g$ проходит пустой ход и затем ведёт планку на $s = T - g = k\,v$. Клавиша сама задаёт значение и сама же выставляет его, как у комптометра.

\begin{center}\small
\begin{tabular}{@{}clll@{}}
\toprule
Планка & Кто ведёт & Зазор & Что делает \\
\midrule
R & ярус продвигающей литеры & 0 & поднимает защёлки, сектора уходят в ноль \\
W & ярус продвигающей литеры & $T - k\,w$ & ставит упор спуска и сектор W \\
B & ярус продвигающей литеры & $T - k\,b$ & ставит сектор B (упор $b$) \\
C & ярус продвигающей литеры & $T - k\,c$ & ставит сектор C (упор $c$) \\
О & ярус мёртвого знака & 0 & отводит ползун до выбранного упора \\
N & ярус подстрочного знака & 0 & вводит упор $c$ вместо упора $b$ \\
\bottomrule
\end{tabular}
\end{center}

\begin{figure}[h]
\centering
\begin{tikzpicture}[x=1mm,y=1mm]
\draw[mech,fill=white] (2,50) rectangle (14,54); \node[lbl,anchor=south] at (8,54.5) {клавиша tinco};
\draw[mech] (14,51) -- (128,51); \fill (128,51) circle (1.2); \node[lbl,anchor=south] at (128,52.5) {ось};
\draw[mech,fill=dead!70] (30,51) rectangle (110,45);
\node[lbl] at (70,48) {кодовая пластина (нижний ярус)};
\def\T{12}\def\k{0.7}
\foreach \x/\v/\nm/\col in {38/12/R/accent,62/10/W/ink,82/5/B/sledA,102/4/C/sledB} {
  \pgfmathsetmacro\g{\T-\k*\v}
  \ifnum\v=12 \def\g{0}\fi
  \draw[mech,fill=white] (\x-2,45) rectangle (\x+2,{45-20+\g});
  \fill[\col!35,draw=\col] (\x-5,{45-20-2.4}) rectangle (\x+5,{45-20});
  \node[lbl,text=\col] at (\x,{45-20-6}) {планка \nm};
}
\node[lbl,anchor=north,align=center] at (38,{45-20-9}) {сброс\\$g{=}0$};
\node[lbl,anchor=north,align=center] at (62,{45-20-9}) {$w{=}10$\\$g{=}5$ мм};
\node[lbl,anchor=north,align=center] at (82,{45-20-9}) {$b{=}5$\\$g{=}8{,}5$ мм};
\node[lbl,anchor=north,align=center] at (102,{45-20-9}) {$c{=}4$\\$g{=}9{,}2$ мм};
\draw[dim] (116,51) -- node[lbl,right,align=left]{$T = 12$ мм\\полный ход} (116,39);
\end{tikzpicture}
\caption{Нижний ярус пластины клавиши tinco. Сброс идёт без зазора и потому раньше замера; чем больше значение, тем меньше зазор. Ступень $k = 0{,}7$~мм на единицу и ход 12~мм — ориентировочные.}
\label{fig:plate}
\end{figure}

\paragraph{Два яруса.} Клавиша регистра поднимает челнок на два ряда и сдвигает мерные планки, чтобы под выступы встал верхний ярус пластины. Ярусы независимы: на одной клавише может быть мёртвый знак (выступы О и N) и продвигающая литера (выступы R, W, B, C). Так раскладка ТПМ-1 переходит в механику без изменений.

\paragraph{Точность.} Значения целые, и всё, что ставится по планкам, имеет фиксатор по целым единицам. Фиксатор доводит положение до ближайшей единицы и прощает неточность зазора меньше половины ступени — 0,35~мм при ступени 0,7~мм.

\subsection{Спуск}

Колесо спуска связано через шестерню с рейкой каретки, пружинный двигатель каретки всё время стремится его повернуть. Зуб колеса соответствует одной единице хода каретки; мелкость 0,4~мм даёт передаточное число, а не мелкий зуб.

\begin{figure}[h]
\centering
\begin{tikzpicture}[x=1mm,y=1mm]
\def\R{22}
\draw[mech,fill=gray!6] (0,0) circle (\R);
\foreach \i in {0,...,39} {\draw[ink!70] ({90+\i*9}:\R) -- ({90+\i*9+4.5}:{\R+2.2}) -- ({90+(\i+1)*9}:\R);}
\fill (0,0) circle (1.2);
\node[lbl] at (0,-5) {колесо спуска};
\draw[arr] (-8,10) arc (130:60:12) node[lbl,above right]{пружина каретки};
\draw[mech,fill=white] (-1.5,{\R+2.4}) -- (1.5,{\R+2.4}) -- (0,{\R-0.5}) -- cycle;
\draw[mech] (0,{\R+2.4}) -- (0,{\R+12}); \node[lbl,anchor=south] at (0,{\R+12}) {держащая собачка Д};
\def\ang{0}
\draw[thinmech,ink!40] (0,0) -- ({\ang}:{\R+8});
\draw[mech,fill=dead] ({\ang}:{\R-0.5}) -- ({\ang+3}:{\R+3}) -- ({\ang-3}:{\R+3}) -- cycle;
\node[lbl,anchor=west] at ({\ang}:{\R+5}) {переставной упор У};
\draw[{Latex[length=1.8mm]}-{Latex[length=1.8mm]},thick,draw=accent] (90:{\R+5}) arc (90:{\ang+4}:{\R+5});
\node[lbl,text=accent] at (45:{\R+10}) {$w$ зубьев};
\draw[mech,fill=gray!15] (-42,-14) circle (5); \node[lbl,anchor=north] at (-42,-20) {шестерня (на валу колеса)};
\draw[mech,dashed,ink!50] (0,0) -- (-42,-14);
\draw[mech] (-64,-8.6) -- (-22,-8.6);
\foreach \x in {-64,-62,...,-24} {\draw[ink!60] (\x,-8.6) -- (\x+1,-6.8);}
\node[lbl,anchor=south] at (-50,-6.5) {рейка каретки};
\node[lbl,anchor=north west,align=left,text width=56mm] at (30,-12) {
\textbf{При отпускании клавиши:}\\
1. Д выходит из зуба;\\
2. колесо идёт до упора У ($w$ зубьев);\\
3. Д падает в зуб;\\
4. У поднимается и возвращается в ноль.};
\end{tikzpicture}
\caption{Спуск. Планка W при нажатии отводит упор У на $w$ зубьев вперёд; при отпускании каретка проходит ровно $w$ единиц. Мёртвые знаки упор не трогают, и каретка стоит.}
\label{fig:esc}
\end{figure}

\subsection{Память: три сектора}

Сектора B, C и W поворачивают планки B, C и W. На дуге каждого сектора зубчатый фиксатор с положениями от 0 до максимального значения и защёлка. Каждый сектор держит свой упор: B — упор отвода для надстрочных знаков, C — для подстрочных, W — упор возврата «назад на литеру».

\begin{figure}[h]
\centering
\begin{tikzpicture}[x=1mm,y=1mm]
\fill (0,0) circle (1.2);
\draw[mech,fill=sledA!10] (0,0) -- (60:30) arc (60:120:30) -- cycle;
\foreach \i in {0,...,9} {\pgfmathsetmacro\a{114-\i*6}
  \draw[ink!70] (\a:30) -- (\a:32);
  \node[lbl,font=\sffamily\tiny] at (\a:35) {\i};}
\draw[mech,fill=white] (0,40) -- (2,32.5) -- (-2,32.5) -- cycle;
\draw[mech] (0,40) -- (0,46); \node[lbl,anchor=south] at (0,46) {защёлка};
\draw[mech] (0,0) -- (-30,-8);
\node[lbl,anchor=east] at (-30,-8) {от планки B};
\draw[mech,sledA,thick] (0,0) -- (40,-6) -- (70,-6);
\draw[mech,fill=sledA!60] (70,-9) rectangle (73,-3);
\node[lbl,anchor=west,text=sledA] at (74,-6) {упор $b$ ползуна};
\draw[arr] (-20,14) arc (150:125:22); \node[lbl,anchor=east] at (-20,14) {пружина к нулю};
\node[lbl,anchor=north west,align=left,text width=72mm] at (-35,-14) {
\textbf{На каждой продвигающей литере:}\\
1. планка R поднимает защёлку, пружина возвращает сектор в ноль;\\
2. планка B ведёт сектор на $b$ делений;\\
3. защёлка падает в ближайшую впадину и держит сектор после отпускания.\\
\textbf{На мёртвом знаке} сектор не трогают: он только служит упором.};
\end{tikzpicture}
\caption{Сектор памяти B. Сектора C и W устроены так же; их ведут планки C и W.}
\label{fig:mem}
\end{figure}

\subsection{Отвод печатающего узла}

При мёртвом знаке планка О через пружинную тягу толкает ползун назад. Ползун останавливается на упоре $b$, а если ярус знака несёт выступ N, — на упоре $c$. Тяга пружинная, поэтому положение задаёт упор, а не глубина нажатия: палец лишь должен дожать клавишу до конца. Затем челнок поворачивается, молоток бьёт, а при отпускании пружина возвращает ползун на нулевой упор.

Так в механизме, задающем положение знака, ровно два звена: нулевой упор и упор сектора. Нулевой упор регулируется один раз и одинаково переносит посадку всех знаков; упоры секторов доводятся фиксатором до целой единицы.

\subsection{Служебные клавиши}

\begin{itemize}
  \item \textbf{Пробелы 4u и 6u} — постоянные пластины: $w = 4$ или $6$, память в ноль.
  \item \textbf{Забой} — обычный: колесо спуска назад на один зуб (1\,u). Память не трогает.
  \item \textbf{Назад на литеру} — храповой толкатель поворачивает колесо спуска назад до упора сектора W, то есть на ширину последней литеры, и сбрасывает всю память. Если после этого ударить ту же литеру через корректирующую ленту, она ляжет точно на прежнее место, а память восстановится.
  \item \textbf{$\pm u$} — качающийся рычажок: пока он отклонён, нулевой упор ползуна смещён на единицу в ту или другую сторону. Ручная поправка для редких случаев (точки под alda, нестандартные сочетания).
\end{itemize}

\subsection{Блокировки}

\begin{itemize}
  \item \textbf{Контроль полного хода.} Начатая клавиша должна дойти до конца: храповая гребёнка не даёт ей вернуться с полпути. Иначе память и упор спуска встанут на неполное значение, а молоток не спустится.
  \item \textbf{Очерёдность.} Пока клавиша не отпущена — подача не прошла, ползун не вернулся, защёлки не легли, — остальные клавиши заперты.
\end{itemize}

%=====================================================================
\section{Ход клавиши}

\begin{figure}[h]
\centering
\begin{tikzpicture}[x=1.05mm,y=1mm]
\draw[ink!50,->] (0,0) -- (118,0);
\foreach \x/\l in {0/0\%,20/,40/,60/,80/,100/100\%} {\draw[ink!50] (\x,0) -- (\x,-1.5) node[lbl,below]{\l};}
\node[lbl,anchor=north] at (50,-5) {ход вниз};
\node[lbl,anchor=north] at (110,-5) {отпускание};
\def\row#1#2{\node[lbl,anchor=east] at (-2,#1) {#2};}
\node[lbl,anchor=west,font=\sffamily\footnotesize\bfseries] at (0,74) {мёртвый знак};
\row{68}{отвод ползуна до упора} \fill[sledA!40] (4,66) rectangle (40,70);
\row{61}{поворот челнока} \fill[accent!30] (20,59) rectangle (70,63);
\row{54}{взвод и спуск молотка} \fill[axisred!50] (60,52) rectangle (100,56); \fill[axisred] (98,52) rectangle (100,56);
\row{47}{возврат ползуна} \fill[sledA!25] (102,45) rectangle (114,49);
\node[lbl,anchor=west,font=\sffamily\footnotesize\bfseries] at (0,40) {продвигающая литера};
\row{34}{сброс памяти (R)} \fill[accent!35] (2,32) rectangle (12,36);
\row{27}{замер W, B, C} \fill[sledB!40] (12,25) rectangle (56,29);
\row{20}{поворот челнока} \fill[accent!30] (20,18) rectangle (70,22);
\row{13}{взвод и спуск молотка} \fill[axisred!50] (60,11) rectangle (100,15); \fill[axisred] (98,11) rectangle (100,15);
\row{6}{защёлки; подача $w$} \fill[sledB!25] (56,4) rectangle (102,8); \fill[axisblue!45] (102,4) rectangle (114,8);
\end{tikzpicture}
\caption{Порядок действий по ходу клавиши (ориентировочно). Удар — тёмная риска в конце хода. Каретка движется только после удара продвигающей литеры; ползун — только при мёртвом знаке.}
\label{fig:stroke}
\end{figure}

Ход клавиши длиннее обычного: его первая часть уходит на замер или отвод. Это та же цена, что у комптометра, где клавиши приводили механизм.

%=====================================================================
\section{Литеры}

\begin{figure}[h]
\centering
\setlength\tgem{26mm}
\begin{tikzpicture}[x=\tgu,y=\tgu]
\input{gen/fig_slugs_adv.tex}
\end{tikzpicture}
\caption{Продвигающие литеры. Красная линия — ось молотка (левый край клетки), серая — ось техты. Под литерой её код.}
\end{figure}

\begin{figure}[h]
\centering
\setlength\tgem{24mm}
\begin{tikzpicture}[x=\tgu,y=\tgu]
\input{gen/fig_slugs_dead.tex}
\end{tikzpicture}
\caption{Мёртвые знаки. Якорь знака — на оси молотка; под знаком указано, на какую величину уходит узел. Черта удвоения для lambe поднята внутрь изгиба.}
\end{figure}

Набор знаков тот же, что в ТПМ-1: один комплект гласных на все тенгвы, черта удвоения и носовая черта в узком и широком вариантах, особая черта для lambe и alda, литеры «ára + гласная».

%=====================================================================
\section{Раскладка}

Раскладка совпадает с ТПМ-1. Мёртвые знаки окрашены по тому, куда уходит узел.

\begin{figure}[h]
\centering
\begin{tikzpicture}[x=1mm,y=1mm]
\input{gen/keyboard.tex}
\end{tikzpicture}
\caption{Раскладка ТПМ-2 (ред. 0.2). Зелёные клавиши — надстрочные знаки (отвод $b$), коричневые — подстрочные (отвод $c$); цвет клавиши — по нижнему регистру. Вверху фиолетовым — верхний регистр.}
\label{fig:kbd}
\end{figure}

%=====================================================================
\section{Набор: два примера}
\label{sec:examples}

В таблицах: «удар» — точка бумаги, куда пришёлся удар (для литеры — левый край клетки, для знака — ось); «отвод» — смещение узла; «каретка» — положение после клавиши; «память» — сектора B / C / W.

\subsection{Квенья: \emph{Elen síla lúmenn' omentielvo}}

\begin{figure}[h]
\centering
\setlength\tgem{12.5mm}
\begin{tikzpicture}[x=\tgu,y=\tgu]
\input{gen/fig_trace_q.tex}
\end{tikzpicture}
\caption{Квенийский режим. Риски под оттиском: красные — удары литер, зелёные и коричневые — удары знаков.}
\end{figure}

{\small\setlength\tgem{5mm}
\begin{longtable}{@{}rlrrrc@{}}
\toprule
\# & Клавиша & Удар & Отвод & Каретка & Память $b/c/w$ \\
\midrule\endhead
\tabinput{gen/tab_trace_q.tex}
\bottomrule
\end{longtable}}

\subsection{Синдарин: \emph{pedo mellon a minno}}

\begin{figure}[h]
\centering
\setlength\tgem{14mm}
\begin{tikzpicture}[x=\tgu,y=\tgu]
\input{gen/fig_trace_s.tex}
\end{tikzpicture}
\caption{Синдаринский режим. В «mellon» черта для lambe бьёт с отводом $c = 5$.}
\end{figure}

{\small\setlength\tgem{5mm}
\begin{longtable}{@{}rlrrrc@{}}
\toprule
\# & Клавиша & Удар & Отвод & Каретка & Память $b/c/w$ \\
\midrule\endhead
\tabinput{gen/tab_trace_s.tex}
\bottomrule
\end{longtable}}

%=====================================================================
\section{Точность}

\subsection{От чего зависит посадка знака}

\begin{center}\small
\begin{tabularx}{\linewidth}{@{}>{\raggedright}p{42mm}>{\raggedright}p{40mm}X@{}}
\toprule
Звено & Чем задано & Погрешность \\
\midrule
рисунок знака на челноке & отливка по якорю шрифта & заводская, одна для всех экземпляров челнока \\
позиция челнока & упор поворота, как у Хаммонда & общая для всех литер: один литероноситель \\
подача каретки & зуб колеса спуска & целые единицы; ошибки не накапливаются в пределах строки \\
отвод узла & упор сектора с фиксатором & целые единицы; неточность зазора до 0,35~мм поглощается фиксатором \\
нулевое положение узла & один регулируемый упор & одна регулировка на всю машинку \\
квантование кода & округление до $u$ & промежуток $\le 0{,}5\,u$ на сторону ($\approx$ 0,2 мм); нижняя ось $\le 0{,}4\,u$ \\
\bottomrule
\end{tabularx}
\end{center}

\subsection{Сравнение с рычажным исполнением}

В рычажном исполнении (ред. 0.1) у каждого из 44 рычагов своё выравнивание, у двух салазок — свои направляющие литер, и между ними и общим сегментом нужен зазор. Ошибка любого из этих звеньев видна только на своих знаках, и её приходится искать по отдельности. На челноке все литеры бьют об одну наковальню одним молотком, и все знаки переносит один ползун; регулировка посадки — одна.

%=====================================================================
\section{Эксплуатация}

\paragraph{Правила для наборщика.}
\begin{enumerate}
  \item Сначала тенгва, потом её знаки — в любом режиме.
  \item Клавишу дожимать до конца; следующую нажимать после отпускания предыдущей (машинка сама не даст иначе).
  \item После пробела знак ложится на край клетки — так ставятся отдельные знаки и пунктуация.
  \item Цифры бьются младшим разрядом вперёд, как пишет рука.
\end{enumerate}

\paragraph{Исправление.} «Назад на литеру», корректирующая лента, тот же удар — литера ложится на прежнее место, и память восстанавливается. Забытую техту к последней букве можно добить в любой момент до следующей буквы.

\paragraph{Смена шрифта и режима.} Челнок снимается и ставится без инструмента, как у Хаммонда. Для режима, где o и u меняются местами, достаточно челнока с переставленными знаками и наклейки на клавиши; механизм кода от этого не меняется, если ширины литер те же.

\paragraph{Обслуживание.} Регулировок три: нулевой упор ползуна (посадка всех знаков), сила пружины молотка (плотность оттиска), натяжение пружины каретки. Фиксаторы секторов и шаг колеса спуска — заводские.

%=====================================================================
\section{Сравнение исполнений}

\begin{center}\small
\begin{tabularx}{\linewidth}{@{}>{\raggedright}p{32mm}>{\raggedright}X>{\raggedright}X>{\raggedright\arraybackslash}X@{}}
\toprule
& ТПМ-1 (электромеханика) & ТПМ-2, ред. 0.1 (рычаги, салазки) & ТПМ-2, ред. 0.2 (челнок) \\
\midrule
сила & мотор & палец & палец, пружина молотка \\
оттиск & ровный (мотор) & зависит от пальца & ровный (пружина молотка) \\
к оси идёт & каретка & салазки знаков & весь печатающий узел \\
память & регистр защёлок & два сектора & три сектора \\
выравнивание & 44 рычага & 44 рычага и двое салазок & один челнок \\
регулировок посадки & по рычагам & по рычагам и салазкам & одна \\
ограничение раскладки & нет & клавиша либо мёртвая, либо продвигающая & нет \\
смена шрифта & нет & нет & сменный челнок \\
исправление & одна литера & $w$ ударов забоя & одна литера \\
скорость & высокая & ниже & ниже рычажной \\
\bottomrule
\end{tabularx}
\end{center}

%=====================================================================
\section{Риски и открытые вопросы}

\begingroup\sloppy
\begin{enumerate}[label=Р\arabic*.]
  \item \textbf{Размер челнока.} Ряд из 22 позиций по 19\,u — около 166~мм развёртки при кегле em~=~1/4~дюйма. Это крупный челнок; масса и время поворота растут вместе с ним. Варианты: кегль 1/6~дюйма (развёртка около 110~мм) или шесть рядов вместо четырёх.
  \item \textbf{Ход и сила клавиши.} Замер (или отвод), поворот челнока и взвод молотка — всё от пальца. Нужен расчёт сил; вероятно, пружины, уравновешивающие мерные планки.
  \item \textbf{Ползун под молотком.} Рама с молотком должна при отводе не перекашиваться и выдерживать удар без сдвига; ползун во время удара стоит на упоре, прижатый пружинной тягой. Нужна проверка жёсткости.
  \item \textbf{Двенадцать уровней подачи.} Varityper обходился тремя; здесь упор спуска ходит на 4--16 зубьев.
  \item \textbf{Скорость.} Механизм Хаммонда и блокировка очерёдности делают машинку медленнее рычажной. Для эльфийского набора это приемлемая цена.
  \item \textbf{Заправка листа.} Молоток стоит за листом, поэтому заправка неудобнее, чем на рычажной машинке.
  \item \textbf{Память на одну литеру.} Знаки ставятся только к последней букве; добить техту к букве двумя раньше можно только забоем и рычажком $\pm u$.
\end{enumerate}
\endgroup

%=====================================================================
\appendix
\section{Коды литер}

$w$ — подача, $a$ — от края клетки до оси техты (для справки), $b$ — отвод для надстрочных знаков, $d$ — смещение нижней оси, $c = b - d$ — отвод для подстрочных. Всё в $u = \mathrm{em}/16$.

{\small\setlength\tgem{5.5mm}
\begin{longtable}{@{}cllrrrrr@{}}
\toprule
& Литера & U+ & $w$ & $a$ & $b$ & $d$ & $c$ \\
\midrule\endhead
\tabinput{gen/tab_codes.tex}
\bottomrule
\end{longtable}}

\section{Отклонённые варианты}

\begin{description}[style=nextline,leftmargin=0pt]
  \item[Рычажная машинка с салазками знаков (ред. 0.1).] Мёртвые рычаги стоят на двух подвижных участках сегмента, которые предыдущая буква отводит на $b$ и $c$. Проще по приводу (нет поворота челнока), но выравнивание приходится держать по 44 рычагам и двум салазкам, сила удара зависит от пальца, а раскладка требует, чтобы клавиша была целиком мёртвой или продвигающей.
  \item[Без нижней оси.] Один сектор памяти, подстрочные знаки по верхней оси. Ошибка подстрочного знака — $d$: 1\,u ($\approx$ 0,4~мм) у основных тенгв, 3\,u у hyarmen. Не проходит по точности.
  \item[Без «назад на литеру».] На один сектор меньше, но исправление требует до 16 ударов забоя. Не проходит по простоте эксплуатации.
  \item[Моноширинная переделка серийной машинки.] Самая простая и тоже точная: у всех литер $b$ одинаково, надстрочные знаки — обычные мёртвые клавиши, отлитые со смещением. Нарушает требование пропорциональных ширин; годится для первого макета.
\end{description}

\section*{Источники}
\begin{itemize}\small
  \item G.\,C.~Mares. History of the Typewriter. London, 1909, с.~143 (устройство Хаммонда: челнок, наковальня, молоток на пружине; цель — выравнивание и равномерный оттиск) — \url{https://tekniskkulturarv.dk/book/21c52e37-c94c-4fd7-8f2f-df7031f6b69b?page=151}.
  \item Antique Typewriters: Hammond 1 (постоянная сила молотка, сменный челнок) — \url{https://antiquetypewriters.com/typewriter/hammond-1-typewriter}.
  \item Collectors Weekly: Beautiful Machine: 1881 Hammond (заправка листа) — \url{https://collectorsweekly.com/articles/?p=16805}.
  \item US~780102. Printing mechanism for type-writers (Хаммонд: молоток, челнок, U-образная рама) — \url{https://patents.google.com/patent/US780102A/en}.
  \item US~2581200. Typewriter carriage spacing control (Varityper) — \url{https://patents.google.com/patent/US2581200}.
  \item US~2904159. Proportional spacing mechanism for typewriters — \url{https://patents.google.com/patent/US2904159A/en}.
  \item HP Museum: Mechanical adders (комптометр) — \url{https://hpmuseum.org/adddir.htm}.
  \item Остальные источники (Alcarin Tengwar, раскладка Смита, IBM Executive) — в документе ТПМ-1.
\end{itemize}

\end{document}
